In this section, we cover the most enduring economic applications of SE and RL. We note that SE has a long history in influencing public policy, while RL has only recently begun to see deployment in business and economic applications.  

\subsection{Economic Applications of Structural Econometrics}

Structural econometric models have influenced policy decisions across various economic domains. In macroeconomic policy, central banks employ structural models as core analytical tools. The Federal Reserve's FRB/US model and the European Central Bank's New Area-Wide Model provide frameworks for monetary policy analysis through counterfactual simulations \citep{Fischer2017, Constancio2018}. As former Federal Reserve Vice Chair Stanley Fischer observed, these models serve as "key tools" for scenario analysis at major central banks \citep{Fischer2017}. Similarly, fiscal policy has incorporated structural approaches through dynamic scoring of tax legislation. The analysis of the 2017 Tax Cuts and Jobs Act employed overlapping-generations models to quantify macroeconomic feedback effects, projecting modest GDP increases and additional tax revenue \citep{TCJA2018}.

Antitrust enforcement has been transformed by structural merger simulation, providing quantitative evidence for regulatory decisions. In U.S. v. H\&R Block (2011), the court explicitly referenced the Justice Department's structural model forecasts in blocking the acquisition of TaxACT \citep{ArnoldPorter2012}. The European Commission has similarly utilized structural models in telecommunications merger reviews, including the Hutchison/Telefonica UK case, where simulations projected potential price increases that influenced regulatory outcomes \citep{DLAPiper2020, BakerBotts2020}.

Market design represents a direct application of structural economics, with theoretical models creating new institutions. The FCC's spectrum auctions exemplify this approach, with auction mechanisms designed to efficiently reallocate radio spectrum \citep{FCC2020}. Educational assignment systems have implemented matching algorithms based on the Gale-Shapley deferred acceptance mechanism, substantially reducing the number of unmatched students when adopted in New York City schools. Similar mechanisms have been employed in healthcare for kidney exchange programs, increasing transplantation rates through multi-patient matching chains.

Social policy decisions often rely on structural models to predict behavioral responses to program changes. The Affordable Care Act's design was informed by microsimulation models projecting coverage expansion and budget impacts \citep{ACA2011}. Labor market policies, including minimum wage regulations, have been evaluated using structural models that estimate employment effects and income distribution changes \citep{NPR2025}. These quantitative predictions provide policymakers with evidence-based projections of reform consequences, though the accuracy of such forecasts remains subject to ongoing empirical validation.

\subsection{Economic Applications of Reinforcement Learning}

Reinforcement learning techniques are beginning to be applied to economic decisions at scale, replacing rule-based or simpler learning methods. 

In transportation, DiDi implemented multi-agent RL for ride-hailing dispatch optimization. Their system, deployed across multiple Chinese cities, achieved $0.5$-$2\%$ improvements in key metrics including driver income and order completion rates compared to previous production baselines \citep{DiDi2019}. At DiDi's scale, these seemingly modest gains translate to substantial efficiency improvements.

Microsoft applied RL to advertising content personalization, notably for Xbox homepage promotions. Their system dynamically selected which promotional content to display to each user based on contextual information. In live A/B testing, the RL-driven approach yielded a $60\%$ higher click-through rate compared to baseline policies \citep{Microsoft2022}, demonstrating the effectiveness of sequential decision-making in user engagement optimization.

In finance, JPMorgan Chase developed LOXM, an RL-driven algorithm for executing large equity orders. Trained on billions of historical trades, the system optimizes execution timing and sizing to minimize market impact while maximizing price efficiency. When deployed on trading desks, LOXM achieved approximately $15\%$ better execution efficiency compared to traditional strategies \citep{JPMorgan2018}, representing one of the first production applications of deep RL in institutional finance.

For e-commerce, Alibaba leveraged RL for dynamic pricing during major shopping events. During their 2018 "Double 11" Shopping Festival, which generated \$30.8 billion in 24 hours, Alibaba deployed RL agents that forecasted product demand and automatically adjusted prices in real-time \citep{Alibaba2019}. This enabled optimal discount strategies and inventory allocation, contributing to a $27\%$ year-over-year growth in sales.

Marketing applications include attempts to optimize customer interactions. \citet{Liu2022} compared model-free RL policies for coupon targeting against conventional approaches in controlled experiments, reporting modest improvements in customer response. \citet{Misra2019}'s successfully demonstrate the value of incorporating incorporating microeconomic choice theory into scalable dynamic price experimentation. 

Google's application of RL to data center cooling demonstrates impact in operational efficiency. By using deep RL to autonomously control cooling equipment in production data centers, Google reduced cooling energy consumption by up to $40\%$ \citep{Google2018}. The system continually adjusts fans, chillers, and other equipment to maintain safe temperatures while minimizing energy usage, yielding a $15\%$ reduction in overall Power Usage Effectiveness (PUE).